\chapter{Loader、接口与可信边界}\label{ch:loader}

一次验证需要调用者选定的完整 source MCM、完整 target MCM、RPX certificate
和独立 Request。当前 RPX 不携带 source/target hash；验证直接读取、检查和使用这两个实际模型。
调用者有责任把正确文件或字节绑定到这两个参数，并在调用之后按 Request 的含义解释结果。
不可信 certificate 只能提交 proof，不能自行改写调用者请求。

\section{入口、处理顺序与判定}

推荐集成入口为 crate \code{kernel_core} 的 \code{api} 模块；
\code{main.rs} 只解析参数、打印 façade 返回的 JSON 和设置 exit code。
文件入口先以受限 owned read 获得完整字节，之后依次执行：
\begin{enumerate}
\item 对 source 与 target 调用 \code{codec::decode_model}，包括全部模型结构验证。
\item 调用 \code{codec::decode_certificate} 完整解析 RPX；任何嵌套分支均不能略过。
\item 调用 \code{rpa::parse_request} 解析独立请求。
\item 核对 certificate 与 Request 的全部 claims、point claims、bijection claims
及 improvement claims 在数量、顺序与每个字段上精确相等。
\item 用实际模型验证各方向及其 complete Returns；随后验证 bijection injectivity、
point witnesses，最后验证 improvements。所有阶段共享同一次工作预算。
\item 只有全部请求通过后返回 \code{VERIFIED}。解析、导入或装配成功不具有该含义。
\end{enumerate}
模型导入可以事先完成。MPS/LP 不在 \code{verify} 内重复解析，也没有 solver 回调。

\begin{longtable}{p{32mm}r p{94mm}}
\toprule 状态 & exit & 含义 \\\midrule\endfirsthead
\toprule 状态 & exit & 含义 \\\midrule\endhead
\code{VERIFIED} & 0 & 独立 Request 的全部数学合同已由此 kernel 接受。\\
\code{CANONICALIZED} & 0 & MPS/LP 解析、模型检查与 MCM 编码完成。\\
\code{ASSEMBLED} & 0 & strict RPA 解析与 RPX 编码完成。\\
\code{DISASSEMBLED} & 0 & RPX 解码并输出 RPA。\\
\code{INSPECTED} & 0 & MCM 解码并报告规模。\\
\code{EXPORTED} & 0 & MCM 解码并导出规范模型 JSON。\\
\code{INVALID} & 2 & 格式、模型、请求匹配、证明或数学义务不满足要求，亦包括不支持的版本/方言。\\
\code{RESOURCE_LIMIT} & 3 & 已公布预算耗尽；不是验证成功，也不是数学反驳。\\
\code{IO_ERROR} & 4 & 文件读写失败。\\
\bottomrule
\end{longtable}
没有 exit 1 的规定状态，也没有单独的 \code{UNSUPPORTED}。
调用者必须看 status，不能把辅助命令的 exit 0 当成模型关系成立。
\code{INVALID} 也不统一表示“目标数学命题为假”：格式错误和 importer 能力限制都属于该状态。

\section{可信精确模型导入}

\subsection{共同规范化与模型身份}
MPS importer 位于 \code{model.rs}，LP importer 位于 \code{lp.rs}；
两者和 MCM codec 都是可信代码。它们用任意精度整数与精确有理数进行解析及规范化，
不经过浮点数，不调用优化器，不修改原始文件。
每个 row/目标表达式把同一变量系数精确相加，按 coordinate ID 排序并删除精确零。
零系数的变量出现仍声明变量，不能因为最后项取消而删除该变量。
名称、原始声明次序、所有变量/row、bounds、integer flags 和完整目标均保留于规范模型。

所有内部目标均最小化。maximization 使\emph{整个}仿射目标取负，常数也要取负。
缺失 bound side 是真正的无限方向；不能用大有限数代替。
互相矛盾的有限 bounds 是空模型的合法表示，不能为了得到可行输入而“修复”。
integer 的有理 bounds 不自动取整，因为 LP relaxation 的含义也须保留。
模型结构合法不证明模型可行，更不证明 source/target 等价。

\subsection{MPS 的 sections、cards 与排列}
MPS headers 必须从第一列开始，data cards 缩进。次序为 NAME，
可选 OBJSENSE/OBJNAME，ROWS，COLUMNS，可选 RHS、RANGES、BOUNDS，最后必须 ENDATA。
OBJSENSE 与 OBJNAME 各至多一次；single value 可放在 header 或下一张缩进 card。
senses 接受 MIN/MINIMIZE、MAX/MAXIMIZE，默认 MIN；headers、keywords、sense 区分大小写。
NAME 未给名称时使用 \code{MODEL}。

必须恰好有一个 N objective row，E/L/G 是 constraint rows。
OBJNAME 若存在必须指定唯一 N row。constraint 顺序为 ROWS 顺序，除去 objective row；
variable 顺序为首次 COLUMNS 出现顺序，重复 column cards 必须连续。
每个 coefficient card 含一或两组 row/value。重复 coefficient 可相加，
重复 row name、未知 row/variable、非连续 column 声明均拒绝。

本方言接受空白分隔 free MPS，以及名称不含空白的普通 fixed-column cards。
\code{COLUMNS/RHS/RANGES} 的 fixed continuation 只有前十四个字符全为空格时才能省略
column/vector name，而且此前须存在可继续的同类 card。
fixed-format 的 embedded-space 名称拒绝；不能靠对齐猜测缺失必要字段。
第一列 \code{*} 开始注释；独立 token 以 \code{\$} 开始时截断后续 inline comment，
附着在名称中的 \code{\$} 不会截断数据。因此本方言不接受以 \code{\$} 开始的数据名称。
空白行和注释行忽略；ENDATA 后非注释内容拒绝。

RHS、RANGES、BOUNDS 每种至多一个 named vector；不默选第一个向量。
同一 row 的重复 RHS 或 range card 拒绝。缺失 RHS 是精确零。
objective row 的 RHS 值 $b$ 生成 objective constant $-b$，然后才按 sense 归一化；
objective 上不能有 range。

\subsection{MPS bounds、markers 与 ranges}
连续 column 默认 $[0,+\infty)$。\code{'MARKER' 'INTORG'} 至
\code{'MARKER' 'INTEND'} 的 column 是 integer；markers 必须配对且不能嵌套。
标记 integer column 如果\emph{完全没有}显式 bound card，默认 $[0,1]$；
一旦存在任何显式 card，未定义 sides 回到一般 $[0,+\infty)$ defaults。
故单独 LI 表示 general integer，不再有隐式 upper 1；标记 integer 的 PL 也显式移除 upper。

\begin{longtable}{p{18mm}p{128mm}}
\toprule bound card & 精确含义 \\\midrule\endfirsthead
\toprule bound card & 精确含义 \\\midrule\endhead
LO / UP & 设置有限 lower / upper；可以为有理数。\\
FX & 两侧设为同一个有限有理数。\\
FR / MI / PL & 两侧 infinite / lower infinite / upper infinite；不允许 numeric field。\\
LI / UI & 设置 lower / upper 并声明 integer；该 card 的值必须是整数。\\
BV & integer 的 $[0,1]$ 与所有显式 bounds 相交，card 顺序无关；可选数值必须恰好为 1。\\
\bottomrule
\end{longtable}
同一 side 重复定义必须相同；冲突的覆盖拒绝。相反 sides 的矛盾保留。
若显式 upper 为负且没有显式 lower card，MPS 默认 lower 变为 $-\infty$；
upper=0 则仍保留默认 lower=0。显式 lower 永远优先于此默认规则。
BV 则始终和 $[0,1]$ 相交：BV+LO 2 得 $[2,1]$，BV+FR 得 $[0,1]$，
BV+UP $-1$ 得 $[0,-1]$，后者不会被负 upper 默认规则变为可行区间。

row RHS 为 $b$、range 为 $r$ 时：
\begin{longtable}{p{27mm}p{119mm}}
\toprule row 类型 & canonical finite sides \\\midrule\endfirsthead
\toprule row 类型 & canonical finite sides \\\midrule\endhead
L & $b-|r|\le a^\top x\le b$。\\
G & $b\le a^\top x\le b+|r|$。\\
E，$r\ge0$ & $b\le a^\top x\le b+r$。\\
E，$r<0$ & $b+r\le a^\top x\le b$。\\
\bottomrule
\end{longtable}
没有 range 时保留原 equality 或单侧 constraint；range=0 生成 equality。

MPS 明确拒绝多个 N row、多个 vectors、初始 blank-name vectors、未知 section/bound type、
SOS、quadratic/conic/indicator、semicontinuous/semiinteger 及其他未定义扩展。
拒绝表明本 importer 没有接受这种方言，不能据此判断模型间关系。

\section{LP v1 的自足文本规则}

\subsection{sections、名称与 coordinate order}
\code{canonical} 依据输入文件扩展名选择 parser：大小写不敏感的 \code{.lp}
调用 LP importer，其他扩展调用 MPS importer；没有内容 sniffing。
LP 输入必须 ASCII；反斜线开始到该物理行末的注释。
除 line endings、tab 外的 control characters 拒绝。
keywords 不区分大小写，每个 section header 独占一行。

必须先 objective，后 constraints，最后 End。Bounds 可选且位于 constraints 之后、
type sections 之前。Binary、General、Semi 可按任意顺序放在 constraints/Bounds 之后，
各至多一次；Bounds 不能出现在 type section 之后。
\begin{longtable}{p{33mm}p{113mm}}
\toprule section & 接受的 aliases \\\midrule\endfirsthead
\toprule section & 接受的 aliases \\\midrule\endhead
min objective & \code{Minimize, Minimum, Min}\\
max objective & \code{Maximize, Maximum, Max}\\
constraints & \code{Subject To, Such That, St, St., S.t.}\\
bounds & \code{Bounds, Bound}\\
binary & \code{Binary, Binaries, Bin}\\
general integer & \code{General, Generals, Gen}\\
empty semi marker & \code{Semi, Semis, Semi-Continuous}\\
end & \code{End}\\
\bottomrule
\end{longtable}
Semi 只允许空 section；任何成员拒绝。End 后只能有空行或注释。

名称以 ASCII letter 或 underscore 开始；后续可用字母、数字及
\code{_.$\#(),;!?@\{\}~\%\&}。名称区分大小写、最长 1,024 B、不截断。
以 e/E 为首，后面缺省或紧接数字、e/E 的 exponent-ambiguous name 拒绝。
名称不能是 section/type keyword、free、inf、infinity 或保留的未支持扩展 keyword：
sos、lazy、scenario/scenarios、indicators、pwl、multi-objectives。
quoted、square-bracket names 和数字/标点起首名称不支持；名称内括号是字面字符，
不是表达式 grouping。

LP 没有 problem-name field，模型名固定为 \code{LP}。
可选 objective label 只检查语法，不另存为模型字段。
variable 顺序按 objective、rows、bounds、type lists 的首次文本出现决定；
constraint 顺序按输入决定。constraint labels 保留且不能重复；
unnamed rows 自动命名为 \code{_lp_row_0}、\code{_lp_row_1} 等。
bounds/type-only 变量也保留。

\subsection{表达式、续行与关系}
表达式由 signed terms 组成：数值常数，或可选系数后接变量。
首项以后每项都要 \code{+} 或 \code{-}。系数与变量必须有空白，
例如 \code{2 x} 合法、\code{2x} 拒绝；operator 自身无需两侧空白，\code{x-y} 合法。
符号可以与 magnitude 分离；objective 和 named row 的 coefficient/variable tokens
可跨物理行。不能使用乘号、power、quadratic brackets、grouping 或函数调用。
零 objective 必须写 \code{0}，不能省略整个表达式。

若第一条 constraint 带 label，只有新 data line 开头的 \code{label:}
开始新 row，其他行继续当前 row。若第一条 constraint 不带 label，
每个物理 constraint line 必须包含一条完整 unnamed relation；
后续 named row 拒绝。不能由 parser 猜测未标号行是续行还是新 row。

单关系为 \code{affine <= affine}、\code{affine >= affine} 或
\code{affine = affine}。reader 减去右端，再把差的 constant 移到 bound。
例如 $x+1\le y+3$ 变为 $x-y\le2$。
range row 允许 \code{constant <= affine <= constant} 或反向的两个 \code{>=}；
方向必须一致，两 endpoints 必须精确化简成常数。
例如 $1\le x+2\le3$ 变为 $-1\le x\le1$；所有导入的 row Expr constant 为零。
named range row 可跨行；strict inequalities 拒绝。

\subsection{bounds 与 integrality}
continuous/general integer 默认 $[0,+\infty)$；显式 bound 替换指定 side。
LP 的负 upper \emph{不会}移除默认 lower zero，故只有 \code{x <= -1} 时为 $[0,-1]$。
这与前节 MPS convention 不同，文件扩展选择是模型语义的一部分。

每个 Bounds record 占一物理行，只约束一个未缩放 variable。
\begin{lstlisting}
0 <= x <= 3
3 >= x >= -2
x <= 4
-2 <= x
x >= -2
4 >= x
x = 1
1 = x
x free
-inf <= x <= +infinity
\end{lstlisting}
inf/infinity 不区分大小写，仅允许负 infinite lower 或正 infinite upper；
infinite equality、正 infinite lower、负 infinite upper 均拒绝。
同一 side 的重复 bound 必须相同；冲突 override 拒绝。
Binary 是 integer 的 $\{0,1\}$ 与 explicit bounds 的交集，
例如 binary x 加 $2\le x\le7$ 为 $[2,1]$，binary x free 为 $[0,1]$。
General 仅声明整数性，不隐含 upper=1。
type lists 为以空白分隔的名称；重复或冲突 type declarations 拒绝。

LP 拒绝未知 sections/tokens、SOS、indicators、scenarios、quadratic/nonlinear syntax、
非空 Semi、缺失必要 section/End、End 后 noncomment content 和未定义 dialect ambiguities。
有理 fraction、双侧 affine relations 是本 importer 明确提供的扩展；
vendor solver 接受范围不自动成为本 kernel 的输入语义。

\section{精确数值与各层资源边界}

MPS、LP 和 RPA 共用 \code{parse_q}。它接受带符号 integers、integer fractions、
有限 decimal（含 \code{.1}、\code{1.}）及 e/E/d/D scientific notation，
包括 \code{1/-2}；分母零拒绝。text 可以是未约分或负分母拼写，解析后规范化；
binary Q 则必须已经规范，不能自动约分。
NaN、infinity coefficients、hexadecimal floats、数字内部 whitespace 和 digit separators 不接受。
无限 bounds 仅由上述 importer grammar 表示。

\begin{longtable}{p{90mm}p{56mm}}
\toprule 层与资源 & 当前上限 \\\midrule\endfirsthead
\toprule 层与资源 & 当前上限 \\\midrule\endhead
façade/CLI 每个 owned input file & 512 MiB\\
RPA / Request text & 256 MiB；JSON nesting capacity 128\\
MPS / LP text & 512 MiB；每物理行 32,768 B\\
LP lexical tokens & 20,000,000\\
model names / variables / rows & 1,024 B / 2,000,000 / 2,000,000\\
MPS ROWS cards（含 objective） & 2,000,000\\
coefficient entries / final model terms & 20,000,000\\
decimal mantissa 或 fraction component digits & 每部分 4,096\\
decimal exponent digits / absolute exponent & 6 / 4,096\\
effective power of ten & 绝对值至多 4,096\\
scalar text bytes & $2\times4,096+32$\\
exact intermediate numerator / denominator & 各 65,536 bits\\
binary codec 文件 / 单 array / 累计 counted items & 1 GiB / 10,000,000 / 20,000,000\\
binary magnitude / name preflight & 每 magnitude 1,700 B / 1 MiB\\
binary proof depth & 64，root=1\\
默认 verification work / live facts / depth & $10^9$ / $10^6$ / 64\\
\bottomrule
\end{longtable}
限制叠加适用，不能用宽松的一层覆盖严格的一层。
例如 text parser 能构造的 65,536-bit Q 可能超过 MCM 的 1,700-byte magnitude 编码预算。
\code{verify_bytes} 不适用路径读取的 512 MiB cap，但仍适用 codec、Request 和验证预算。
theorem constructors 还有自己的维数和工作限制；生成模型也必须满足同一个 Model validation。
预算耗尽没有浮点 fallback、wraparound 或“尽力接受”的模式。
重要 allocations 使用 \code{try_reserve}；本接口不承诺第三方任意精度库的所有内部
allocation failures 或宿主强制终止都能恢复成 JSON。

work 在各方向、lift、injectivity、point refutation 与 improvements 间累积，
嵌套 scopes 不重置。live facts 限制当前可见事实；proof nesting 与 JSON parser nesting
是两种不同限制。CLI 可以设 work 和 live facts，不能设置 depth。
版本最低值与资源限制也不同：V10 本身不提高默认预算。

\section{真实 CLI 与 Rust façade}

当前命令如下；二进制名为 \code{rpx}，仓库构建命令为
\code{cargo build --release -p kernel-core --bin rpx}。
\begin{lstlisting}
rpx canonical input.mps output.mcm
rpx canonical input.lp output.mcm
rpx assemble input.rpa output.rpx
rpx disassemble input.rpx output.rpa
rpx inspect-model model.mcm
rpx export-model model.mcm output.json
rpx verify --work 1000000000 --live-facts 1000000 \
  source.mcm target.mcm proof.rpx request.json
\end{lstlisting}
verify 的两个 options 可出现在 subcommand 后任何位置；paths 仍须恰好四个，
顺序为 source、target、proof、request。数值分别解析为 u64 和 usize。
其他 command 不接受这些 flags，没有 \code{--profile}、\code{--direction}、
hash 参数或 certificate 自定 budget。未知 option、缺参数、路径数量错误返回 usage INVALID。

façade 的路径函数与 CLI 同名操作对应：
\code{canonicalize_path}、\code{assemble_paths}、\code{disassemble_paths}、
\code{inspect_model_path}、\code{export_model_paths}、\code{verify_paths}。
需要内存集成时使用 \code{api::verify_bytes(source,target,proof,request,limits)}，
前三项为 byte slices、request 为 text；不触及文件系统。
入口返回 \code{Result<serde_json::Value, api::Error>}。
更底层的 \code{verify::verify_request} 接受内存中的 Model/Certificate/Request 与 Limits，
并自行再次验证两个完整模型；
带位置的接口为 \code{verify_request_with_diagnostics}。
legacy \code{verify} 接受 Claim slice，只请求方向，不接受 points/bijections/improvements。
独立 LP contradiction 和完整 injectivity 库入口分别为
\code{verify_lp_refutation}、\code{verify_injectivity}，并有相应 diagnostics 版本。
低层 \code{verify_injectivity} 只检查给定完整 map 的单射性；它不独立建立总可行性、
与证书索引的绑定或完整双射。完整合约验收使用 \code{verify_request}。

路径读取先检查 metadata length，再用限长 read 确保文件增长也不超过 cap；
实际用于验证的是取得的 owned bytes。
这不是操作系统级多文件事务：如需 source/target/proof/request 的一致业务快照，
由调用者先获取并通过 \code{verify_bytes} 绑定。
output 在内容全部生成后创建父目录并写入指定文件；路径函数可能替换现有 output，
不承诺 atomic rename 或自动 backup。

\section{报告、诊断与模型 JSON}

每次 CLI invocation 向 stdout 打印一个 JSON value。verify 成功包括
status、claims/points/bijections counts、statistics、limits、
\code{model_decode_us}、\code{proof_decode_us}、\code{execute_us}、\code{total_us}；
路径调用另有 \code{read_us}，内存调用没有该字段。
非空 improvement 请求增加 \code{improvements} count；新增 V10 statistics 零值省略。
这些 timings 是本次调用诊断量，不能单独解释为隔离 benchmark。

failure 必含 status、message、code；已知时增加 location 与 limits。
\begin{lstlisting}
{"status":"INVALID",
 "message":"INVALID obligation mismatch: expected Eq; obligation 3",
 "code":"obligation_mismatch",
 "location":{"claim":0,"block":0,"obligation":3},
 "limits":{"work":1000000000,"live_facts":1000000,"depth":64}}
\end{lstlisting}
code 是错误词汇生成的 snake case 诊断辅助，不是冻结枚举；
稳定分类应以 status 为准。Location 全部 indices 为零基：
claim/point 是 Request 各自数组下标，block 是顶层 certificate Block 下标，
step 是最内层活动 Block/Derivation 的 Step 下标，
obligation 是当前 terminal 的实际检查次序下标。
没有相应语义的位置省略字段。improvement、fiber refutation 等不是顶层方向 Block，
不能从不存在的 claim/block 字段推断坐标。
diagnostics 只是 best-effort error attribution，不参与数学决定、统计或预算。
CLI 在有效 verify 参数解析后的失败上附加 effective limits；
库调用者可用 \code{Error::with_limits} 附加。

\code{inspect-model} 报告模型名、variables/rows、integer count 和 row/目标总 nonzeros；
\code{export-model} 通过 \code{model::to_json} 输出
name、variables、rows、objective。每个 rational 为 canonical scalar string，
bounds 缺侧为 null，terms 是有序的 \code{[variable,coefficient]} pairs。
JSON object key 顺序不具有语义，数组次序保留模型坐标。
这是 export schema，不能直接使用 Model 的普通 Serialize derive 作为同样的 scalar interchange，
也没有当前 CLI 的 JSON-model import command。

\section{信任结论与外部证据}

接受边界包含精确数值 parser/算术、MPS/LP importers、Model validation、
MCM/RPX codec、Request matching、scalar/scope checker、完整 obligation construction、
point 与 injectivity checker，以及各静态 lift/theorem 的数学 lemma 和 Rust 实现。
扩展的 old-scope premise checking 只证明其前提；生成模型合法且 child proof 通过，
仍不能独立证明联合存在性 lemma。这一项必须明确列入 TCB。

certificate generator、求解器、Python oracle、model-specific structure scripts、
外部 SHA receipts、benchmark 和 corpus replay 都不向 kernel 提供 proof facts。
独立测试、字节攻击、构造 oracle 与 source review 是实现可靠性的证据，
不是 Rust 实现的形式验证。V9/V10 的冻结来自 2026-10-14 的明确项目决定，
其基线与复核记录见 \code{audit/PROTOCOL_V9_V10_FREEZE.md}。
验证已准许请求不自动获得未请求的双向等价、最优性、可行性或 global inclusion。

本章以 \code{api.rs}、\code{main.rs}、\code{model.rs}、\code{lp.rs}、
\code{verify.rs} 为实现依据；输入方言详解对应 \code{docs/MPS_V1.md}、
\code{docs/LP_V1.md}，接口说明对应 \code{docs/API.md}。
